\clearpage
\addcontentsline{toc}{section}{第3セット 採点基準}

\noindent
{\Large\sffamily\bfseries 第3セット\ \ 採点基準}
\par\vspace{2mm}

\hdB{配点一覧（100点満点／200点換算併記）}
\begin{center}
\small
\renewcommand{\arraystretch}{1.3}
\begin{tabular}{@{}clcc@{}}
\toprule
\textbf{大問} & \textbf{分野・主題} & \textbf{小問の配点} & \textbf{計} \\
\midrule
1 & 数学A・B 整数と指数数列（合同式・隣接3項間漸化式・最大公約数・平方数） & $4 + 5 + 5 + 6$ & 20 \\
2 & 数学A・B・III 確率と漸化式（正三角柱・状態縮約・連立確率漸化式と極限） & $4 + 5 + 5 + 6$ & 20 \\
3 & 数学II・C・III 二次曲線と微積分（放物線の準線からの直交接線・極線・面積・回転体） & $4 + 5 + 5 + 6$ & 20 \\
4 & 数学III 空間図形と積分（連立不等式の立体・2方向の断面・最短距離・回転体） & $4 + 5 + 6 + 5$ & 20 \\
5 & 数学III 定積分漸化式と極限・無限級数（はさみうちの原理・交代級数の和） & $4 + 4 + 6 + 6$ & 20 \\
\bottomrule
\end{tabular}
\end{center}

\hdB{自己採点 成績評価基準（志望校判定・偏差値目安・到達水準）}
\par\vspace{1.0mm}\noindent
{\footnotesize 本模試（難関大二次・冠模試レベル／120分完全記述式）の自己採点得点に対応する合格目安と学習指針です。}
\par\vspace{1.5mm}
\begin{center}\footnotesize
\renewcommand{\arraystretch}{1.28}
\begin{tabular}{@{}cclp{26em}@{}}
\toprule
\textbf{得点（100点）} & \textbf{換算点（200点）} & \textbf{判定・偏差値目安} & \textbf{実力評価と今後の学習課題} \\
\midrule
$75 \sim 100$ 点 & $150 \sim 200$ 点 & \textbf{Sランク}（偏差値 70〜） & \textbf{東大・京大・難関医学部 A判定レベル}。第1問の平方数否定・第2問のランダムウォーク・第4問の回転体まで，誘導の少ない高度な難関大実戦問題を論理の破綻なく解き切れています。名大オープン等各大学の個別冠模試でも上位合格が狙えます。 \\
$60 \sim 74$ 点 & $120 \sim 148$ 点 & \textbf{Aランク}（偏差値 63〜69） & \textbf{旧帝大理系・難関医学部 合格水準}。各大問の(1)〜(3)を着実にクリアし，実戦的な方針点を確実に稼ぎ出せています。第1問(4)の背理法・第4問(4)の立体回転など，最後の一歩を詰める記述力を磨きましょう。 \\
$40 \sim 59$ 点 & $80 \sim 118$ 点 & \textbf{Bランク}（偏差値 55〜62） & \textbf{旧帝大理系・難関大二次ボーダーライン}。問題の難しさに圧倒されず，方針点・途中式を部分点として粘り強く拾う記述戦略が奏功しています。時間配分（各大問20〜25分）を見直し，完答大問を2題作る訓練をしましょう。 \\
$25 \sim 39$ 点 & $50 \sim 78$ 点 & \textbf{Cランク}（偏差値 48〜54） & \textbf{発展総合演習への挑戦段階}。各大問の(1)(2)は解けていますが，分野融合の応用段階で手が止まっています。第2セット（標準〜応用編）の完全復習と，過去問・冠模試過去問の各大問の誘導の読み取りを徹底強化しましょう。 \\
$25$ 点未満 & $50$ 点未満 & \textbf{Dランク}（偏差値 〜47） & \textbf{要・第1セットおよび第2セットの徹底反復}。難問に挑む前に，標準的な典型題の論述を減点なく書き切る基礎実戦力を固めることが急務です。 \\
\bottomrule
\end{tabular}
\end{center}

\rubricintro{整数の剰余分類における合同式の根拠（平方数の余りの候補など），連立漸化式における対称性の根拠と初期条件の確認，回転体の軸と断面図形の包含関係（円板か円環か），はさみうちの原理における両側の不等式評価の厳密性など，難関大実戦レベルの記述答案に求められる論理的整合性を厳格に採点する。}

% =====================================================================
% 大問1
% =====================================================================
\hdA{大問1（20点）整数と指数数列・最大公約数・平方数}

\Rh{1}{1}{4}
\Seikai{$a_{1}=9$，$a_{2}=33$，$a_{3}=141$，$3$ の倍数の証明}
\begin{saiten}
\Sp{$a_{1}=9$，$a_{2}=33$，$a_{3}=141$ をすべて正しく求めた}{+2}
\Sp{$5\equiv -1, 2\equiv -1\pmod{3}$ より $a_{n}\equiv (-1)^{n}+(-1)^{n+1}=0\pmod{3}$（または数学的帰納法・因数分解）を示した}{+2}
\Ssep
\Sm{$a_{1}, a_{2}, a_{3}$ のいずれかの計算を誤った}{-1}
\end{saiten}

\Rh{1}{2}{5}
\Seikai{漸化式の成立，$n\equiv 1\pmod{3}$（すなわち $n=3m-2$）}
\begin{saiten}
\Sp{$7a_{n+1}-10a_{n}=a_{n+2}$ の代数計算を正しく完了した}{+2}
\Sp{$a_{n+3}\equiv -a_{n}\pmod{9}$（または二項展開 $a_{n}\equiv 3(-1)^{n}(2n+1)\pmod{9}$）を導いた}{+2}
\Sp{$a_{1}\equiv 0, a_{2}\equiv 6, a_{3}\equiv 6\pmod{9}$ から $n\equiv 1\pmod{3}$ を結論づけた}{+1}
\Ssep
\Sm{$n\equiv 1\pmod{3}$ の十分性のみを示し，必要性の論証がない}{-2}
\end{saiten}

\Rh{1}{3}{5}
\Seikai{$a_{n+1}-2a_{n}=3\cdot 5^{n}$，$5a_{n}-a_{n+1}=3\cdot 2^{n+1}$，最大公約数は $3$}
\begin{saiten}
\Sp{$a_{n+1}-2a_{n}=3\cdot 5^{n}$ および $5a_{n}-a_{n+1}=3\cdot 2^{n+1}$ をともに正しく求めた}{+2}
\Sp{最大公約数 $g$ が $3$ の倍数かつ $3\cdot 5^{n}$ と $3\cdot 2^{n+1}$ の公約数であることを指摘した}{+2}
\Sp{$\gcd(5^{n}, 2^{n+1})=1$ から $g=3$ であることを導いた}{+1}
\Ssep
\Sm{$g=3$ であることを具体的な $n=1,2$ から類推したのみで一般の $n$ について論証していない}{-3}
\end{saiten}

\Rh{1}{4}{6}
\Seikai{$n=1$}
\begin{saiten}
\Sp{$n=1$ で $a_{1}=9=3^{2}$ が平方数であることを確認した}{+1}
\Sp{$n$ が偶数のとき $\bmod 5$（または $\bmod 8$）で平方数の剰余条件を満たさないことを示した}{+2}
\Sp{$n$ が奇数（$n\geqq 3$）のとき $\bmod 8$（または $\bmod 9$）で平方数の剰余条件を満たさないことを示した}{+2}
\Sp{以上より $n=1$ のみであることを結論づけた}{+1}
\Ssep
\Sm{$n$ が偶数の場合または奇数の場合の一方の検討が欠落している}{-3}
\end{saiten}

% =====================================================================
% 大問2
% =====================================================================
\hdA{大問2（20点）正三角柱上のランダムウォークと確率漸化式}

\Rh{2}{1}{4}
\Seikai{$a_{n+1}=\dfrac{2}{3}b_{n}+\dfrac{1}{3}d_{n}$，$b_{n+1}=\dfrac{1}{3}a_{n}+\dfrac{1}{3}b_{n}+\dfrac{1}{3}e_{n}$，$d_{n+1}=\dfrac{1}{3}a_{n}+\dfrac{2}{3}e_{n}$，$e_{n+1}=\dfrac{1}{3}b_{n}+\dfrac{1}{3}d_{n}+\dfrac{1}{3}e_{n}$}
\begin{saiten}
\Sp{対称性 $b_{n}=c_{n}, e_{n}=f_{n}$ を利用して遷移関係を正しく立式した}{+2}
\Sp{4つの推移確率漸化式をすべて正しく導いた}{+2}
\Ssep
\Sm{遷移確率の係数を誤った（各誤りごとに）}{-1}
\end{saiten}

\Rh{2}{2}{5}
\Seikai{$p_{n+1}=\dfrac{1}{3}p_{n}+\dfrac{1}{3}$，$p_{n}=\dfrac{1}{2}+\dfrac{1}{2}\left(\dfrac{1}{3}\right)^{n}$}
\begin{saiten}
\Sp{$p_{n+1}=a_{n+1}+2b_{n+1}$ に代入し，$d_{n}+2e_{n}=1-p_{n}$ を用いて $p_{n+1}=\dfrac{1}{3}p_{n}+\dfrac{1}{3}$ を導いた}{+2}
\Sp{特性方程式を解き，初項 $p_{0}=1$（または $p_{1}=\dfrac{2}{3}$）から一般項を求めた}{+2}
\Sp{$p_{n}=\dfrac{1}{2}+\dfrac{1}{2}\left(\dfrac{1}{3}\right)^{n}$ を正しく求めた}{+1}
\Ssep
\Sm{初項を誤り，公比の指数をずらした}{-2}
\end{saiten}

\Rh{2}{3}{5}
\Seikai{$a_{n}+d_{n}=\dfrac{1}{3}$，$b_{n}+e_{n}=\dfrac{1}{3}$（$n\geqq 1$）}
\begin{saiten}
\Sp{$a_{n+1}+d_{n+1}=\dfrac{1}{3}(a_{n}+d_{n})+\dfrac{2}{3}(b_{n}+e_{n})$ を導いた}{+2}
\Sp{全確率の和 $(a_{n}+d_{n})+2(b_{n}+e_{n})=1$ を代入して右辺が $\dfrac{1}{3}$ になることを示した}{+2}
\Sp{$n\geqq 1$ でともに $\dfrac{1}{3}$ の定数となることを結論づけた}{+1}
\Ssep
\Sm{$b_{n}+e_{n}$ の係数 $2$ を忘れて全確率の和を崩した}{-2}
\end{saiten}

\Rh{2}{4}{6}
\Seikai{$a_{n}, b_{n}, d_{n}, e_{n}$ の一般項，$\lim\limits_{n\to\infty}a_{n}=\dfrac{1}{6}$}
\begin{saiten}
\Sp{$r_{n+1}=a_{n+1}-b_{n+1}$ について $d_{n}-e_{n}=-r_{n}$ を用い，$r_{n+1}=-\dfrac{2}{3}r_{n}$ を導いた}{+2}
\Sp{$r_{n}=\dfrac{1}{2}\left(-\dfrac{2}{3}\right)^{n}$（$n\geqq 1$）を求めた}{+1}
\Sp{$a_{n}=\dfrac{1}{3}(p_{n}+2r_{n})$ 等から $a_{n}, b_{n}, d_{n}, e_{n}$ の一般項をすべて求めた}{+2}
\Sp{極限値 $\lim\limits_{n\to\infty}a_{n}=\dfrac{1}{6}$ を求めた}{+1}
\Ssep
\Sm{$r_{1}$ の初期値を誤った}{-2}
\end{saiten}

% =====================================================================
% 大問3
% =====================================================================
\hdA{大問3（20点）放物線の準線からの直交2接線・面積・回転体}

\Rh{3}{1}{4}
\Seikai{$\alpha+\beta=2a$，$\alpha\beta=-1$，直交性の証明}
\begin{saiten}
\Sp{接線の方程式 $y=tx-\dfrac{1}{2}t^{2}$ に点 $\Pt{P}\left(a, -\dfrac{1}{2}\right)$ を代入し，$t^{2}-2at-1=0$ を導いた}{+2}
\Sp{解と係数の関係から $\alpha+\beta=2a, \alpha\beta=-1$ を求めた}{+1}
\Sp{傾きの積が $\alpha\beta=-1$ であることから直交性を示した}{+1}
\Ssep
\Sm{判別式 $D/4=a^{2}+1>0$ による異なる2解の存在確認がない}{-1}
\end{saiten}

\Rh{3}{2}{5}
\Seikai{直線 $\Pt{QR}: y=ax+\dfrac{1}{2}$，定点 $\Pt{F}\left(0, \dfrac{1}{2}\right)$}
\begin{saiten}
\Sp{$2$ 点 $\Pt{Q}, \Pt{R}$ を通る直線の傾き $\dfrac{1}{2}(\alpha+\beta)=a$ を求めた}{+2}
\Sp{直線の方程式 $y=ax+\dfrac{1}{2}$ を求めた}{+2}
\Sp{$x=0$ で $y=\dfrac{1}{2}$ となり定点 $\Pt{F}\left(0, \dfrac{1}{2}\right)$ を通ることを示した}{+1}
\Ssep
\Sm{定点 $\Pt{F}$ の $y$ 座標の符号を誤った}{-1}
\end{saiten}

\Rh{3}{3}{5}
\Seikai{$S(a)=\dfrac{2}{3}(a^{2}+1)^{3/2}$，最小値 $\dfrac{2}{3}$，$\triangle\Pt{PQR}=1$}
\begin{saiten}
\Sp{$S(a)=\dfrac{1}{12}(\beta-\alpha)^{3}$ と $\beta-\alpha=2\sqrt{a^{2}+1}$ から $S(a)=\dfrac{2}{3}(a^{2}+1)^{3/2}$ を求めた}{+2}
\Sp{$a=0$ のとき最小値 $\dfrac{2}{3}$ をとることを求めた}{+1}
\Sp{$a=0$ のときの $3$ 点の座標から $\triangle\Pt{PQR}=1$ を求めた}{+2}
\Ssep
\Sm{$\beta-\alpha$ の展開計算ミス}{-1}
\end{saiten}

\Rh{3}{4}{6}
\Seikai{$V_{x}=\dfrac{2}{5}\pi$，$V_{y}=\dfrac{\pi}{24}$}
\begin{saiten}
\Sp{$V_{x}=\pi\int_{-1}^{1}\left\{\left(\dfrac{1}{2}\right)^{2}-\left(\dfrac{1}{2}x^{2}\right)^{2}\right\}dx$ を正しく立式し，$V_{x}=\dfrac{2}{5}\pi$ を求めた}{+3}
\Sp{$V_{y}=\pi\int_{0}^{1/2}\left\{\left(y+\dfrac{1}{2}\right)^{2}-2y\right\}dy$ を正しく立式した}{+2}
\Sp{積分を実行して $V_{y}=\dfrac{\pi}{24}$ を求めた}{+1}
\Ssep
\Sm{回転体の $\pi$ の付け忘れ}{-1}
\Sm{$V_{y}$ の積分区間や境界の取り違え}{-2}
\end{saiten}

% =====================================================================
% 大問4
% =====================================================================
\hdA{大問4（20点）不等式で表された空間図形の断面・最短距離・回転体}

\Rh{4}{1}{4}
\Seikai{$S(t)=\dfrac{1}{2}\left(1-t^{2}\right)^{2}$，$V=\dfrac{8}{15}$}
\begin{saiten}
\Sp{切り口が直角二等辺三角形であることから $S(t)=\dfrac{1}{2}(1-t^{2})^{2}$ を求めた}{+2}
\Sp{偶関数の定積分を実行して $V=\dfrac{8}{15}$ を求めた}{+2}
\Ssep
\Sm{積分計算の係数ミス}{-1}
\end{saiten}

\Rh{4}{2}{5}
\Seikai{$T(u)=\dfrac{4}{3}(1-u)^{3/2}$，体積の一致の確認}
\begin{saiten}
\Sp{切り口の境界 $z=(1-u)-x^{2}$ から幅 $2\sqrt{1-u}$ を導き，$T(u)=\dfrac{4}{3}(1-u)^{3/2}$ を求めた}{+3}
\Sp{$\int_{0}^{1}T(u)\,du=\dfrac{8}{15}$ を計算し，(1) の体積と一致することを示した}{+2}
\Ssep
\Sm{$\dfrac{1}{6}$ 公式の係数ミス}{-1}
\end{saiten}

\Rh{4}{3}{6}
\Seikai{$d(t)=\sqrt{\dfrac{t^{4}+1}{2}}$，最小値 $\dfrac{\sqrt{2}}{2}$（点 $\left(0, \dfrac{1}{2}, \dfrac{1}{2}\right)$）}
\begin{saiten}
\Sp{$\Pt{OP}^{2}=t^{2}+\dfrac{(1-t^{2})^{2}}{2}+\dfrac{(y-z)^{2}}{2}$ から $y=z=\dfrac{1-t^{2}}{2}$ で最小となることを示した}{+2}
\Sp{$d(t)=\sqrt{\dfrac{t^{4}+1}{2}}$ を求めた}{+2}
\Sp{$t=0$ で最小値 $\dfrac{\sqrt{2}}{2}$ をとること，および点 $\left(0, \dfrac{1}{2}, \dfrac{1}{2}\right)$ を求めた}{+2}
\Ssep
\Sm{$y=z$ の条件の根拠がない}{-1}
\end{saiten}

\Rh{4}{4}{5}
\Seikai{$R(t)=1-t^{2}$，$W=\dfrac{16}{15}\pi$}
\begin{saiten}
\Sp{切り口の三角形と中心 $\Pt{H}(t, 0, 0)$ の距離の最大値 $R(t)=1-t^{2}$ を求めた}{+2}
\Sp{回転体が半径 $R(t)$ の円板であることを示し，$W=\int_{-1}^{1}\pi\{R(t)\}^{2}dt$ と立式した}{+2}
\Sp{$W=\dfrac{16}{15}\pi$ を求めた}{+1}
\Ssep
\Sm{回転体の $\pi$ の付け忘れ}{-1}
\end{saiten}

% =====================================================================
% 大問5
% =====================================================================
\hdA{大問5（20点）定積分漸化式・はさみうちの原理・無限級数の和}

\Rh{5}{1}{4}
\Seikai{$I_{0}=\log\dfrac{4}{3}$，$I_{1}=1-3\log\dfrac{4}{3}$}
\begin{saiten}
\Sp{$I_{0}=\int_{0}^{1}\dfrac{1}{x+3}dx=\log\dfrac{4}{3}$ を求めた}{+2}
\Sp{$I_{1}=\int_{0}^{1}\left(1-\dfrac{3}{x+3}\right)dx=1-3\log\dfrac{4}{3}$ を求めた}{+2}
\Ssep
\Sm{対数の底または真数の計算ミス}{-1}
\end{saiten}

\Rh{5}{2}{4}
\Seikai{$I_{n+1}+3I_{n}=\dfrac{1}{n+1}$}
\begin{saiten}
\Sp{$I_{n+1}+3I_{n}=\int_{0}^{1}\dfrac{x^{n}(x+3)}{x+3}dx=\int_{0}^{1}x^{n}dx$ を示した}{+2}
\Sp{積分を実行して $\dfrac{1}{n+1}$ を導いた}{+2}
\Ssep
\Sm{積分の原始関数の係数ミス}{-1}
\end{saiten}

\Rh{5}{3}{6}
\Seikai{不等式の証明，$\lim\limits_{n\to\infty}I_{n}=0$，$\lim\limits_{n\to\infty}nI_{n}=\dfrac{1}{4}$}
\begin{saiten}
\Sp{$0\leqq x\leqq 1$ で $\dfrac{1}{4}\leqq\dfrac{1}{x+3}\leqq\dfrac{1}{3}$ より不等式を示し，はさみうちで $\lim\limits_{n\to\infty}I_{n}=0$ を求めた}{+2}
\Sp{$I_{n+1}\leqq I_{n}$ と (2) から $\dfrac{1}{4(n+1)}\leqq I_{n}$ および $I_{n}\leqq\dfrac{1}{4n}$ を導いた}{+2}
\Sp{はさみうちの原理により $\lim\limits_{n\to\infty}nI_{n}=\dfrac{1}{4}$ を求めた}{+2}
\Ssep
\Sm{$I_{n+1}\leqq I_{n}$ の根拠（$x^{n+1}\leqq x^{n}$）がない}{-1}
\end{saiten}

\Rh{5}{4}{6}
\Seikai{等式の証明，無限級数の和 $\log\dfrac{4}{3}$}
\begin{saiten}
\Sp{両辺に $\dfrac{(-1)^{k-1}}{3^{k}}$ を掛けて和をとり，望遠鏡和で $I_{0}+\dfrac{(-1)^{n-1}}{3^{n}}I_{n}$ となることを示した}{+3}
\Sp{$\lim\limits_{n\to\infty}\dfrac{(-1)^{n-1}}{3^{n}}I_{n}=0$ を明記した}{+1}
\Sp{無限級数の和が $I_{0}=\log\dfrac{4}{3}$ となることを結論づけた}{+2}
\Ssep
\Sm{剰余項の極限が $0$ になる根拠を省いた}{-2}
\end{saiten}

